
\section{A conjectured law for the number of steps}\label{sec:countlaw}

Theorem~\ref{thm:C} fixes the order of growth from below. The exact number of steps is described by a law that is conjectured and
tested by predictions recorded before the computations. In the coordinate $t=2\sqrt x$ the Poisson noise has nearly unit variance,
so the range $[0,A]$ has Fisher length $2\sqrt A$ and half-length $\sqrt A$. The levels of the optimal input are spaced like those
of the Gaussian channel with unit noise and input in $[-a,a]$, here with $a=\sqrt A$~\cite{Paper1}. The exception is the level at silence: there the output
is degenerate, and the levels next to it form a boundary layer. Let $A_g(K)$ be the half-length at which the Gaussian count changes
from $K-1$ to $K$. Let $A_K$ be the peak count at which the Poisson count does. The law reads
\begin{equation}
\sqrt{A_K}=A_g(K)+c(K),\qquad c(K)=\tfrac12\,E\bigl(g_{\rm eff}(K)\bigr)+0.002,\qquad g_{\rm eff}(K)=A_g(K+1)-A_g(K-1).\label{eq:countlaw}
\end{equation}
In words, the $K$-th step of the neuron appears where the Gaussian channel gains its $K$-th level, measured in the Fisher
half-length, shifted by a small wall constant $c(K)$ that accounts for the boundary layer at silence. To predict $A_K$ one takes
$A_g(K-1)$, $A_g(K)$ and $A_g(K+1)$ from the Gaussian channel, forms $g_{\rm eff}(K)$, evaluates $c(K)$ and squares
$A_g(K)+c(K)$. Here $E(g)$ is an excess defined from two half-line chains, small models of the levels next to one end of the range.
Each chain has thirty free levels next to the wall and frozen levels at spacing $g$ beyond them, one chain with Poisson noise and one
with Gaussian noise. The excess is the difference
between their last free levels, measured in the Fisher coordinate. We call this the finite-chain variant. The fit $E(g)=0.9840-0.2347\,g$
to the chain's own excess on $1.43\le g\le2.35$ never sees a Poisson count change, so the rows of Table~\ref{tab:countlaw} are
predictions out of sample. The term $0.002$ is a correction for the finite length of the chain, estimated
from chains of 20, 30 and 40 levels. A second variant, the infinite-chain variant, uses the fit for the limit of an infinitely long chain, $E(g)=1.0049-0.2454\,g$ without the
correction. As a check after the fact, over $K=5$ to $24$ the term $\frac12E(g_{\rm eff})$ alone matches the measured $c(K)$, which runs from $0.237$ to $0.335$, with a root-mean-square deviation of $0.0015$, about half a percent.
How the count changes are computed is described in Section~\ref{sec:methods}.

\begin{table}[tp]
\caption{Count changes of the Poisson channel predicted by the law~\eqref{eq:countlaw} before they were computed. $A_K$ is
the peak count at which the optimal input gains its $K$-th mass point, without spontaneous count. The tolerance of $0.15$ in $A$ was fixed with the first prediction, before
any of these changes was computed; since a change $\delta$ in $A$ moves $\sqrt A$ by about $\delta/(2\sqrt A)$, it is $0.15/(2\sqrt A)$
in the half-length, $0.0037$ at $K=25$ and $0.0017$ at $K=67$, close to the scatter of the law; relative to $A$ it is $3.6\times10^{-4}$
at $K=25$ and $7.5\times10^{-5}$ at $K=67$. Row $21$ was predicted from an earlier local fit of the same law and row $22$ from the law without the correction
$0.002$; the later rows use the finite-chain variant as written.}\label{tab:countlaw}
\centering\small
\begin{tabular}{rrrrl}
\hline
$K$ & predicted $A_K$ & computed $A_K$ & difference & outcome\\
\hline
21 & 307.99 & 308.0874 & $+0.0974$ & held\\
22 & 333.14 & 333.1959 & $+0.0559$ & held\\
23 & 358.97 & 358.9655 & $-0.0045$ & held\\
24 & 385.39 & 385.3810 & $-0.0090$ & held\\
25 & 412.43 & 412.4271 & $-0.0029$ & held\\
47 & 1145.508 & 1145.6157 & $+0.1077$ & held\\
67 & 2001.381 & 2001.5878 & $+0.2068$ & failed\\
\hline
\end{tabular}
\end{table}

Table~\ref{tab:countlaw} lists the seven predictions made so far. Six held. The prediction for the 67th level failed: the
change lies $0.207$ above the predicted value, outside the tolerance. Its miss and that at $K=47$ grow with $K$, and both lie where
$g_{\rm eff}$ is $1.04$ to $1.14$, below the range on which the chain fit was made. The miss at $K=67$ is $1.0\times10^{-4}$ of $A$.
The infinite-chain variant lies $0.05$ above the measured changes at $K=47$ and $K=67$, a nearly constant offset, while the misses of the finite-chain
variant grow as $g_{\rm eff}$ falls. Read after the failure, this points to the finite-chain correction, estimated from chains of 20 to 40
levels, carried outside its range; it does not single out the form of the law. The law is a conjecture that holds to $0.15$ in $A$
up to $47$ levels and fails at $67$. The computed count changes themselves are being turned into proved statements. At the lower
end of the bracket of the 25th level, $A=412.42711264916346$, a computer-assisted proof shows that the optimal
input has exactly 24 mass points (Section~\ref{sec:certmeth}). Three further
changes were predicted with both variants before any computation, the finite-chain value first and the infinite-chain value second:
$A_{26}$ at $440.085$ and $440.141$, $A_{48}$ at $1184.393$ and
$1184.557$, and $A_{68}$ at $2048.241$ and $2048.503$. They will decide between the variants and are not yet computed. These peak counts lie far above those of real fibres, whose
staircases have only a few steps. The predictions test the mathematics of the law, while the fibres of Section~\ref{sec:an} test its
consequences for neurons.

Through the Fisher length the law carries the growth of the Gaussian count to the neuron. The Gaussian count grows like $a^{4/3}$ at
half-length $a$ (proved from below and measured in~\cite{Paper1}), and with $a=\sqrt A$ this becomes a Poisson count of order $A^{2/3}$, the order that Theorem~\ref{thm:C} proves from below.
